\documentclass[11pt]{amsart}
\usepackage{amsfonts, amssymb, amsmath}
\usepackage{mathrsfs}
\usepackage{tikz-cd}
\usepackage{url}
\usepackage{bbm}
\usepackage[hidelinks]{hyperref}
\usepackage{adjustbox}
\newcommand{\spa}{\medskip}
\usepackage[utf8]{inputenc}

%\usepackage{showkeys}

\linespread{1.05}
\tikzset{
	labl/.style={anchor=south, rotate=90, inner sep=.5mm}
}

\linespread{1.05}
\setlength\parindent{0pt}
\usepackage[left=25mm,right=25mm,top=38mm,bottom=34mm]{geometry}



%Rings & co
\newcommand{\C}{\mathbb{C}}
\newcommand{\F}{\mathbb{F}}
\newcommand{\Fp}{\mathbb{F}_p}
\newcommand{\Fq}{\mathbb{F}_q}
\newcommand{\Fqd}{\mathbb{F}_{q^d}}
\newcommand{\Fqn}{\mathbb{F}_{q^n}}
\newcommand{\Fqs}{\mathbb{F}_{q^s}}
\newcommand{\Fqsp}{\mathbb{F}_{q^{t}}}
\newcommand{\Ga}{\mathbb{G}_a}
\newcommand{\LL}{\mathbb{L}}
\newcommand{\KK}{\mathbb{K}}
\newcommand{\Q}{\mathbb{Q}}
\newcommand{\Z}{\mathbb{Z}}
\newcommand{\N}{\mathbb{N}}
\newcommand{\R}{\mathbb{R}}
\newcommand{\A}{\mathbb{A}}
\newcommand{\El}{E_\lambda}
\newcommand{\Elx}{\El^{(x_0)}}
\newcommand{\hZ}{\widehat{\Z}}
\newcommand{\Zl}{\mathbb{Z}_{\ell}}
\newcommand{\Ql}{\mathbb{Q}_{\ell}}
\newcommand{\Kbar}{\bar{K}}
\newcommand{\kbar}{{\bar{k}}}
\newcommand{\HH}{\mathrm{H}}
\newcommand{\Qlbar}{\bar{\mathbb{Q}}_{\ell}}
\newcommand{\Qpbar}{\bar{\mathbb{Q}}_{p}}
\newcommand{\Qlp}{\mathbb{Q}_{\ell'}}
\newcommand{\Qq}{\mathbb{Q}_{q}}
\newcommand{\Zq}{\mathbb{Z}_{q}}
\newcommand{\Zp}{{\mathbb{Z}_{p}}}
\newcommand{\Qp}{{\mathbb{Q}_{p}}}
\newcommand{\Zmpn}{\mathbb{Z}/p^n}
\newcommand{\Qqn}{\mathbb{Q}_{q^n}}
\newcommand{\Qqs}{\mathbb{Q}_{q^s}}
\newcommand{\Qqsp}{\mathbb{Q}_{q^{t}}}
\newcommand{\Qlpbar}{\bar{\mathbb{Q}}_{\ell'}}
\newcommand{\Qbar}{\bar{\mathbb{Q}}}
\newcommand{\Gm}{\mathbb{G}_m}
\newcommand{\G}[1]{\mathbb{G}_{m,#1}}
\newcommand{\sym}{{\mathrm{sym}}}

%mathcals
\newcommand{\calA}{\mathcal{A}}
\newcommand{\calB}{\mathcal{B}}
\newcommand{\calC}{\mathcal{C}}
\newcommand{\calE}{\mathcal{E}}
\newcommand{\calF}{\mathcal{F}}
\newcommand{\calG}{\mathcal{G}}
\newcommand{\calH}{\mathcal{H}}
\newcommand{\calL}{\mathcal{L}}
\newcommand{\calO}{\mathcal{O}}
\newcommand{\calP}{\mathcal{P}}
\newcommand{\calM}{\mathcal{M}}
\newcommand{\calN}{\mathcal{N}}
\newcommand{\calX}{\mathcal{X}}
\newcommand{\calY}{\mathcal{Y}}
\newcommand{\calU}{\mathcal{U}}
\newcommand{\calV}{\mathcal{V}}
\newcommand{\ucalV}{\underline{\mathcal{V}}}
\newcommand{\ucalVz}{\underline{\mathcal{V}_0}}
\newcommand{\calW}{\mathcal{W}}
\newcommand{\calS}{\mathcal{S}}

%categories
\newcommand{\Weil}{\mathbf{Weil}}
\newcommand{\VVec}{\mathbf{Vec}}
\newcommand{\bfC}{\mathbf{C}}
\newcommand{\Rep}{\mathbf{Rep}}
\newcommand{\Fisoc}{\mathbf{F\textrm{-}Isoc}}
\newcommand{\Isoc}{\mathbf{Isoc}}
\newcommand{\LS}{\mathbf{LS}}
\newcommand{\Coef}{\mathbf{Coef}}
\newcommand{\Foi}{\Fisoc^\textrm{\dag}}
\newcommand{\oi}{\Isoc^\textrm{\dag}}
\newcommand{\myim}{\mathrm{im}}


%DeclareMathOperator

\DeclareMathOperator{\gsms}{gss}
\DeclareMathOperator{\sms}{ss}
\DeclareMathOperator{\hol}{hol}
\DeclareMathOperator{\smooth}{smooth}
\DeclareMathOperator{\Spec}{Spec}
\DeclareMathOperator{\Gal}{Gal}
\DeclareMathOperator{\Hom}{Hom}
\DeclareMathOperator{\Ext}{Ext}
\DeclareMathOperator{\Ker}{Ker}
\DeclareMathOperator{\Isom}{Isom}
\DeclareMathOperator{\End}{End}
\DeclareMathOperator{\Aut}{Aut}
\DeclareMathOperator{\Mod}{Mod}
\DeclareMathOperator{\GL}{GL}
\DeclareMathOperator{\id}{id}
\DeclareMathOperator{\Image}{Im}
\DeclareMathOperator{\ord}{ord}
\DeclareMathOperator{\Tr}{Tr}
\DeclareMathOperator{\Sym}{Sym}
\DeclareMathOperator{\DGal}{DGal}
\DeclareMathOperator{\rk}{rk}
\DeclareMathOperator{\res}{res}
\DeclareMathOperator{\cores}{cores}





%others
\newcommand{\EEnd}{\mathcal{E}nd}
\newcommand{\HHom}{\mathcal{H}om}
\newcommand{\et}{\mathrm{\acute{e}t}}
\newcommand{\cst}{\mathrm{cst}}
\newcommand{\type}{\mathrm{type}}
\newcommand{\crys}{{\mathrm{crys}}}
\newcommand{\Pic}{\mathrm{Pic}}
\newcommand{\Br}{\mathrm{Br}}
\newcommand{\NS}{\mathrm{NS}}
\newcommand{\inv}{{\mathrm{inv}}}
\newcommand{\fppf}{\mathrm{fppf}}
\newcommand{\pil}{\pi_1^\lambda}
\newcommand{\pie}{\pi_1^{\et}}
\newcommand{\iso}{\xrightarrow{\sim}}
\newcommand{\kSch}{k\textrm{-}\mathrm{Sch}}
\newcommand{\Mor}{\mathrm{Mor}}
\newcommand{\ks}{{k_s}}
\newcommand{\ki}{{k_i}}
\newcommand{\kperf}{{k^{\mathrm{perf}}}}
\newcommand{\mupn}{{\mu_{p^n}}}
\newcommand{\kinv}{{k\textrm{-}\inv}}
\newcommand{\Tp}{\mathrm{T}_p}
\newcommand{\Vp}{\mathrm{V}_p}
\newcommand{\Tl}{\mathrm{T_\ell}}
\newcommand{\Tpet}{\mathrm{T^{\et}_p}}
\newcommand{\pdiv}{{}}
\begin{document}

\newtheorem{theo}[subsection]{Theorem}
\newtheorem*{theo*}{Theorem}
\newtheorem{conj}[subsection]{Conjecture}

\newtheorem{prop}[subsection]{Proposition}
\newtheorem{lemm}[subsection]{Lemma}
\newtheorem*{lemm*}{Lemma}
\newtheorem{coro}[subsection]{Corollary}


\theoremstyle{definition}
\newtheorem{defi}[subsection]{Definition}
\newtheorem*{defi*}{Definition}
\newtheorem{hypo}[subsection]{Hypothesis}
\newtheorem{rema}[subsection]{Remark}
\newtheorem{exam}[subsection]{Example}
\newtheorem{nota}[subsection]{Notation}
\newtheorem{cons}[subsection]{Construction}
\numberwithin{equation}{subsection}

\setcounter{tocdepth}{1}


\title{Boundedness of the $p$-primary torsion of the Brauer group of an abelian variety}
\author{Marco D'Addezio}
\date{\today}

\address{Institut de Mathématiques de Jussieu-Paris Rive Gauche, SU - 4 place Jussieu, Case 247, 75252 Paris}
\email{daddezio@imj-prg.fr}

\maketitle

\begin{abstract}
We prove that the $p^\infty$-torsion of the transcendental Brauer group of an abelian variety over a finitely generated field of characteristic $p>0$ is bounded. This answers a question asked by Skorobogatov and Zarhin for abelian varieties. To do this, we prove a flat variant of the Tate conjecture for divisors. In the text, we also study other geometric Galois-invariant $p^\infty$-torsion classes of the Brauer group which are not in the transcendental Brauer group. These classes, in contrast with our main theorem, can be infinitely $p$-divisible. We explain how the existence of these $p$-divisible towers is naturally related to the failure of surjectivity of the specialisation morphisms of the Néron--Severi group in characteristic $p$.
\end{abstract}
	\tableofcontents
\section{Introduction}

In this article we want to study problems related to the finiteness of the $p$-primary torsion of the Brauer group of abelian varieties in positive characteristic $p$. If $k$ is a finite field and $A$ is an abelian variety over $k$ it is well-known that the Brauer group of $A$, defined as $\Br(A):=H^2_\et(A,\Gm)$, is a finite group, \cite[Prop. 4.3]{Tat94}. The main input for this result is the Tate conjecture for divisors proved by Tate in \cite{Tat66}. If $k$ is replaced by a finitely generated field extension of $\Fp$ one can not expect instead $\Br(A)$ to be finite. Indeed, when $k$ is infinite, $\Br(k)$ is an infinite group and it injects into $\Br(A)$. Even $\Br(A)/\Br(k)$ might be infinite since it contains a subgroup isomorphic to $H^1_\et(k,\Pic_{A/k})$. On the other hand, if $\Br(A_{\ks})^{k}$ is the image of the natural morphism $\Br(A)\to \Br(A_{\ks})$, where $\ks$ is a separable closure of $k$ in an algebraic closure $\kbar$, the group $\Br(A_{\ks})^{k}[\tfrac{1}{p}]$ is finite by \cite[Thm. 16.2.3]{CS21}. The group $\Br(A_{\ks})^{k}$ is also known as the \textit{transcendental Brauer group} of $A$. In general, the $p$-torsion of $\Br(A_{\ks})^{k}$ is not finite. Nonetheless, we prove the following finiteness result.

\begin{theo}[Theorem \ref{fini-Brau:t}]\label{i-main:t}
Let $A$ be an abelian variety over a finitely generated field $k$ of characteristic $p>0$. The transcendental Brauer group $\Br(A_{\ks})^{k}$ has finite exponent. In particular, if the Witt vector cohomology group $H^2(A_{\kbar},W\calO_{A_\kbar})$ is a finite $W(\kbar)$-module, then $\Br(A_{\ks})^{k}$ is finite.
\end{theo}

The condition on $H^2(A_{\kbar},W\calO_{A_\kbar})$ is a condition to remove “supersingular pathologies”. It is equivalent to the condition that the \textit{formal Brauer group} of $A_\kbar$ is of finite height (cf. \cite{AM77}). This is satisfied, for example, when the $p$-rank of $A$ is $g$ or $g-1$, where $g$ is the dimension of $A$, \cite[Cor. 6.3.16]{Ill83}. As a consequence of Theorem \ref{i-main:t}, we deduce that subgroup of Galois fixed points, denoted by $\Br(A_{\ks})^{\Gamma_k}$, has finite exponent as well (Corollary \ref{galo-fixe:c}). This answers \cite[Ques. 1 and Ques. 2]{SZ08} for abelian varieties. 

\spa

It is very natural to ask whether the same is true for $\Br(A_{\kbar})^{\Gamma_k}$. In this case, we provide a counterexample which is built using the following result.

\begin{prop}[Proposition \ref{counter:p}]\label{i-counter:p}
	Let $B$ be an abelian variety over $k$. Write $A$ for $B\times B$ and $\Tp(\Br(A_\kbar))$ for the $p$-adic Tate module of $\Br(A_\kbar)$. There is a natural exact sequence		$$0\to \Hom(B,B^\vee)_\Zp\to \Hom(B_\kbar[p^\infty],B_\kbar^\vee[p^\infty])^{\Gamma_k}\to \Tp(\Br(A_\kbar))^{\Gamma_k}.$$
	

\end{prop}

	The proposition implies, for example, that when $\End(B)=\Z$, the $\Gamma_k$-module $\Tp(\mathrm{Br}(A_{\bar{k}}))$ admits non-zero Galois fixed points (Corollary \ref{count:c}). This provides a counterexample to the conjecture in \cite[§7.3.1]{Ulm14}. Note that if we replace $\Tp(\mathrm{Br}(A_{\bar{k}}))$ with the $\ell$-adic Tate module $\Tl(\mathrm{Br}(A_{\bar{k}}))$, where $\ell$ is a prime different from $p$, then $\Tl(\mathrm{Br}(A_{\ks}))=\Tl(\mathrm{Br}(A_{\bar{k}}))$ has no non-trivial Galois fixed points instead. 
	
	\spa
	
	
	These “exceptional classes” in $\Tp(\Br(A_\kbar))^{\Gamma_k}$ are naturally related to the specialisation morphism of the Néron--Severi group. We recall the following theorem, which was proved in \cite[Thm. 5.2]{And96} and \cite{MP12} in characteristic $0$ and in \cite{Amb18} and \cite{Chr18} in positive characteristic. 
	
	\begin{theo}[Andr\'e, Ambrosi, Christensen]\label{i-spec:t}
Let $K$ be an algebraically closed field which is not an algebraic extension of a finite field, $X$ a finite type $K$-scheme, and $\calY\to X$ a smooth and proper morphism. For every geometric point $\bar{\eta}$ of $X$ there is an $x\in X(K)$ such that $\rk_\Z(\NS(\calY_{\bar{\eta}}))=\rk_\Z(\NS(\calY_x))$.
	\end{theo}
	
	As it is well-known, the theorem is false when $K=\bar{\F}_p$ (see \cite[Rmk. 1.12]{MP12}). What we prove is that, in the known counterexamples, the elements in $\Tp(\Br(A_\kbar))^{\Gamma_k}$ explain the failure of Theorem \ref{i-spec:t}. More precisely, we prove the following result.

\begin{theo}[Theorem \ref{NS:t}]\label{i-NS:t}
	Let $X$ be a connected normal scheme of finite type over $\Fp$ with generic point $\eta=\Spec(k)$ and let $f:\calA\to X$ be an abelian scheme over $X$ with constant slopes. For every closed point $x=\Spec(\kappa)$ of $X$ we have $$\rk_\Z(\NS(\calA_{\bar{x}})^{\Gamma_\kappa})-\rk_\Z(\NS(\calA_{\bar{\eta}})^{\Gamma_k})\geq \rk_\Zp ( \Tp(\Br(\calA_{\bar{\eta}}))^{\Gamma_k}).$$
\end{theo}
Note that in the inequality, the left term is “motivic”, while the right term comes from some $p$-adic object which, as far as we know, has no $\ell$-adic analogue. 

\spa

To prove these results we use a flat version of the Tate conjecture. For every $n$, let $H^2_\fppf(A_\kbar,\mupn)^k$ be the image of the extension of scalars morphism $H^2_\fppf(A,\mupn)\to H^2_\fppf(A_\kbar,\mupn)$.

\begin{theo}[Theorem \ref{fppf-tate:t}]\label{i-fppf-tate:t}
	Let $A$ be an abelian variety over a finitely generated field $k$ of characteristic $p$. The cycle class map $$\NS(A)_\Zp\xrightarrow{c_1} \varprojlim_n H^2(A_\kbar,\mupn)^k$$ is an isomorphism. In addition, $\Tp(\Br(A_\kbar)^k)=0$.
\end{theo}


We obtain this result by using the crystalline Tate conjecture for abelian varieties proved by de Jong, \cite[Thm. 2.6]{deJ98}. The main issue that we have to overcome is the lack of a good comparison between crystalline and fppf cohomology of $\Zp(1)$ over imperfect fields. To avoid this problem, we exploit the fact that we are working with abelian varieties. In this special case, the comparison is constructed using the $p$-divisible group of $A$ (and its dual). 

\spa

Note also that in the $\ell$-adic setting, in order to prove that $\Br(A_\ks)^{k}[\ell^\infty]$ is finite knowing the $\ell$-adic Tate conjecture, one needs to show that the Galois representation on $H^2_\et(A_\kbar,\Ql(1))$ is semi-simple. This is used to preserve the exactness of the sequence $$0\to \NS(A_\kbar)_{\Ql}\to H^2_\et(A_\kbar,\Ql(1))\to \Tl(\mathrm{Br}(A_{\bar{k}}))[\tfrac{1}{\ell}]\to 0$$when taking Galois fixed points. In our case, using the groups $H^2_{\fppf}(A_\kbar,\mupn)^k$, we overcome this problem. The technical issue that we have to solve instead is that it is not clear a priori whether $H^2_{\fppf}(A,\Zp(1))\to \varprojlim_n H^2_{\fppf}(A_\kbar,\mupn)^k$ is surjective. We solve this issue in Proposition \ref{tricky:p}, where we reduce to the case when $A$ is the Jacobian of a curve. This was inspired by the proof of \cite[Thm. 2.1]{CS13}.

\subsection{Outline of the article}In §\ref{prel-resu:s} we prove some general results on the cohomology of fppf sheaves. In particular, we prove Corollary \ref{Br-inj:c}, which is a first result on the relation between the Brauer group of a scheme over $\ks$ and $\kbar$. In this section, we also prove in Proposition \ref{Kummer-k-inv:p} the exactness of some fundamental sequences for the groups $H^2(X_\kbar,\mupn)^k$. In §\ref{cons-morp:s}, we construct a morphism which relates $H^2(A,\Zp(1))$ with $\Hom(A[p^\infty],A^\vee[p^\infty])$ and we prove basic properties as Proposition \ref{comm-squa:p} and Proposition \ref{imag:p}. In §\ref{main-resu:s}, we prove the flat variant of the Tate conjecture (Theorem \ref{fppf-tate:t}) and the finiteness result for the transcendental Brauer group (Theorem \ref{fini-Brau:t}). Finally in §\ref{spec:s}, we look at the relation of our results with the theory of specialisation of Néron--Severi groups. In particular, we prove Theorem \ref{NS:t}.
 

\subsection{Acknowledgments} I thank Emiliano Ambrosi for the discussions we had during the writing of \cite{AD18} which inspired this article, Matthew Morrow and Kay Rülling for many enlightening conversations on the cohomology of $\Z_p(1)$, and Ofer Gabber, Luc Illusie, Peter Scholze, and Takashi Suzuki for answering some questions on the fppf site. I also thank Jean-Louis Colliot-Thélène, Takashi Suzuki, and Alexei Skorobogatov for very useful comments on a first draft of this article.

	\spa

The author was funded by the Deutsche Forschungsgemeinschaft (EXC-2046/1, project ID: 390685689 and DA-2534/1-1, project ID: 461915680) and by the Max-Planck Institute for Mathematics.
\section{Notation}

If $k$ is a field, we write $\kbar$ for a fixed algebraic closure of $k$ and $\ks$ (resp. $\ki$) for the separable (resp. purely inseparable) closure of $k$ in $\kbar$. We denote by $\Gamma_k$ the absolute Galois group of $k$. If $x$ is a $k$-point of a scheme, we denote by $\bar{x}$ the induced $\kbar$-point. For an abelian group $M$, we write $\Tp(M)$ for the $p$-adic Tate module of $M$, which is the projective limit $\varprojlim_{n}M[p^n]$, we write $\Vp(M)$ for $\Tp(M)[\tfrac{1}{p}]$, and we write $M^\wedge$ for the $p$-adic completion of $M$. If $M$ is endowed with a $\Gamma_k$-action, we denote by $M^{\Gamma_k}$ the subgroup of fixed points. For a scheme $X$ and an fppf sheaf $\calF$, we denote by $H^{\bullet}(X,\calF)$ the fppf cohomology groups. When $X=\Spec(k)$ we simply write $H^{\bullet}(k,\calF)$. If $f:X\to Y$ is a morphism of schemes we denote by $R^{\bullet}f_*\calF$ the higher direct image functors from $(\mathrm{Sch}/X)_\fppf\to (\mathrm{Sch}/Y)_\fppf$. Finally, we write $X^{\mathrm{perf}}$ for the projective limit $\varprojlim(...\xrightarrow{F}X\xrightarrow{F}X\xrightarrow{F}X)$ where $F$ is the absolute Frobenius of $X$.
\section{Preliminary results}\label{prel-resu:s}

In this section we start with some results that we will use later on. We suppose that $k$ is a field of positive characteristic $p$ and $q:X\to \Spec(k)$ is a scheme over $k$.


\begin{lemm}\label{fppf-exte-scal:l}
Let $\calF$ be a sheaf over $(\mathrm{Sch}/k)_\fppf$ such that $q_*\calF_{X}=\calF$ and suppose that $X$ has a $k$-rational point. The group  $H^0(k,R^1q_*\calF_{X})$ is canonically isomorphic to $H^1(X,\calF_X)/H^1(k,\calF)$. In addition, the natural morphism $H^2(X,\calF_X)\to H^0(k,R^2q_*\calF_{X})$ sits in an exact sequence 
$$0\to K\to H^2(X,\calF_X)\to H^0(k,R^2q_*\calF_{X})\to H^2(k,R^1q_*\calF_{X})$$
where $K$ is an extension of $H^1(k,R^1q_*\calF_{X})$ by $H^2(k,\calF)$. 
\end{lemm}

\begin{proof}
We consider the Leray spectral sequence
$$E^{i,j}_2=H^i(k,R^jq_{*}\calF_X)\Rightarrow H^{i+j}(X,\calF_X).$$ 	The morphisms $E^{i,0}_2=H^i(k,q_*\calF_X)=H^i(k,\calF)\to H^i(X,\calF_X)$ are injective since $X$ admits a $k$-rational point. We deduce that $E^{1,1}_2=E^{1,1}_\infty$ and $E^{2,0}_2=E^{2,0}_\infty$. This implies that the kernel of $H^2(X,\calF_X)\to E^{0,2}_\infty$ is an extension of $E^{1,1}_2$ by $E^{2,0}_2$, as we wanted. The obstruction for the map $H^2(X,\calF_X)\to E^{0,2}_2=H^0(k,R^2q_*\calF_{X})$ to be surjective lies in $E^{2,1}_2=H^2(k,R^1q_*\calF_X)$. This concludes the proof.
\end{proof}
\begin{lemm}\label{limi:l}
	Let $G$ be a commutative finite type group scheme over $k$. If $X$ is quasi-compact quasi-separated, then for every $i$ there is a natural injective morphism $H^0(k,R^i q_*G_X)\hookrightarrow H^{i}(X_\kbar,G_{X_\kbar})$.
\end{lemm}
\begin{proof}
	We first note that since $G$ is of finite type over $k$, for every $k$-algebra $R$ we have that $$G({X_R})=\mathrm{colim}_{R'\subseteq R\  \mathrm{fin. gen.} }G(X_{R'}).$$ Since $X$ is quasi-compact quasi-separated, by \cite[Tag 073E]{SP22} this implies that $$H^i(X_R,G_{X_R})=\mathrm{colim}_{R'\subseteq R\  \mathrm{fin. gen.} }H^i(X_{R'},G_{X_{R'}}).$$
	Let $\calF$ be the presheaf of abelian groups which associates to $T\in (\mathrm{Sch}/k)_{\fppf}$ the group $H^i(X_T,G_T)$. We write $\calF^+$ for the presheaf which sends $T$ to $\mathrm{colim}_{\calU\in \mathrm{Cov}(T)}\breve{H}^0(\calU,\calF)$ where $\mathrm{Cov}(T)$ is the set of fppf covers of $T$. When $T=\Spec(k')$, the finite field extensions $k''/k'$ provide a cofinal set of covers of $T$. We deduce that $$H^0(k',\calF^+)= \mathrm{colim}_{k''/k'\  \mathrm{fin.}}\mathrm{eq} (H^i(X_{k''},G)\rightrightarrows H^i(X_{k''\otimes_{k'} k''},G)).$$ Combining the previous observation with the fact that filtered colimits of abelian groups commute with finite limits, we get $$H^0(k',\calF^+)=\mathrm{eq} (H^i(X_{\kbar},G)\rightrightarrows H^i(X_{\kbar\otimes_{k'} \kbar},G_{X_{\kbar\otimes_{k'} \kbar}}))\subseteq H^i(X_{\kbar},G_{X_\kbar}).$$ As $R^i q_*G$ is the sheafification of $\calF$, we have that $$H^0(k,R^i q_*G_X)=\mathrm{colim}_{{k'/k\  \mathrm{fin.}}}\breve{H}^0(k',\calF^+)\subseteq H^{i}(X_\kbar,G_{X_\kbar}).$$ This ends the proof.\end{proof}

\begin{coro}\label{Br-inj:c}Suppose $k$ separably closed. If $X$ is a proper connected scheme with a $k$-rational point such that $\Pic_{X/k}$ is smooth, then the morphism $\Br(X)\to \Br(X_{\kbar})$ is injective.
\end{coro}

\begin{proof}
Thanks to Lemma \ref{limi:l} it is enough to prove that the morphism $\Br(X)\to H^0(k,R^2q_*\Gm)$ is injective. By Lemma \ref{fppf-exte-scal:l} applied to $X$ and $\Gm$ we deduce that the kernel of $\Br(X)\to H^0(k,R^2q_*\Gm)$ is an extension of $H^1(k,\Pic_{X/{k}})$ by $\Br(k)$. Since $k$ is separably close, both groups vanish, the first one because we assumed $\Pic_{X/{k}}$ smooth. This shows that $\Br(X)\to H^0(k,R^2q_*\Gm)$ is an injection, as we wanted.
\end{proof}

\begin{cons}
		The Kummer exact sequences for $X$ and $X_{\kbar}$ induce the following commutative diagram with exact rows
	\begin{equation}\label{comp:e}\begin{tikzcd}
			0\arrow[r] &\Pic(X)/p^n \arrow[d,""]\arrow[r] & H^2_{}(X,\mupn)\arrow[d,""]\arrow[r] & \mathrm{Br}(X)[p^n]\arrow[r]\arrow[d,""] & 0\\
			0\arrow[r]& \Pic(X_\kbar)/p^n \arrow[r]& H^2_{}(X_\kbar,\mupn)\arrow[r]& \mathrm{Br}(X_\kbar)[p^n]\arrow[r]& 0.\
		\end{tikzcd}
	\end{equation}

We write 	\begin{equation}\label{1:e}0\to(\Pic(X_\kbar)/p^n)^k\to H^2(X_\kbar,\mupn)^{k}\to (\Br(X_{\kbar})[p^n])^{k}\to 0.	\end{equation}
for the sequence obtained by taking images of the vertical arrows. Note that a priori $(\Br(X_{\kbar})[p^n])^{k}$ is smaller than $\Br(X_\kbar)^k[p^n]$, where $\Br(X_\kbar)^k:=\myim(\Br(X)\to \Br(X_\kbar))$.

\spa

We also consider the following commutative diagram with exact rows
\begin{equation*}\label{diag:e}\begin{tikzcd}
		0\arrow[r] &\Pic(X)^\wedge \arrow[d,""]\arrow[r] & H^2_{}(X,\Zp(1))\arrow[d,""]\arrow[r] & \Tp(\mathrm{Br}(X))\arrow[r]\arrow[d,""] & 0\\
		0\arrow[r]& \Pic(X_\kbar)^\wedge \arrow[r]& H^2_{}(X_\kbar,\Zp(1))\arrow[r]& \Tp(\mathrm{Br}(X_\kbar))\arrow[r]& 0\
	\end{tikzcd}
\end{equation*}

obtained by taking the limit of the diagrams (\ref{comp:e}) for various $n$. We denote by 	\begin{equation}\label{2:e}
	0\to(\Pic(X_\kbar)^\wedge)^k\to H^2(X_\kbar,\Z_p(1))^{k}\to \Tp(\Br(X_{\kbar}))^{k}\to 0.
\end{equation}
the sequence obtained by taking images of the vertical arrows.
\end{cons}




\begin{prop}\label{Kummer-k-inv:p} Let $X$ be a proper connected scheme over $k$ with a $k$-rational point. The sequences (\ref{1:e}) and (\ref{2:e}) are exact.
\end{prop}
\begin{proof}


If $K_{1,n}$ is the kernel of $H^2_{}(X,\mupn)\to H^2(X_\kbar,\mupn)$ and $K_{2}$ is the kernel of $\Br(X)\to \Br(X_\kbar)$, to prove that (\ref{1:e}) is exact we have to show that $K_{1,n}\to K_{2}[p^n]$ is surjective. By Lemma \ref{fppf-exte-scal:l}, we have the following commutative diagram with exact rows 

	\begin{equation*}\begin{tikzcd}
		0\arrow[r] &\Br(k)[p^n]\arrow[d,""]\arrow[r] &  K_{1,n}\arrow[d,""]\arrow[r] & H^1_{}(k,\Pic_{X/k}[p^n])\arrow[r]\arrow[d,""] & 0\\
		0\arrow[r]& \Br(k) \arrow[r]& K_2\arrow[r]& H^1_{}(k,\Pic_{X/k})\arrow[r]  & 0.\\
	\end{tikzcd}
\end{equation*}

Since the image of $H^1(k,\Pic_{X/k}[p^n])\to H^1_{}(k,\Pic_{X/k})$ is $H^1_{}(k,\Pic_{X/k})[p^n]$, the morphism of exact sequences factors through $$0\to \Br(k)[p^n]\to K_2[p^n]\to H^1_{}(k,\Pic_{X/k})[p^n]\to 0$$ and $K_{1,n}\to K_2[p^n]$ is surjective.

\spa

We look now at (\ref{2:e}). The kernel of $H^2(X,\Zp(1))\to H^2(X_\kbar,\Zp(1))$ is $\varprojlim_n K_{1,n}$ and the kernel of $\Tp(\Br(X))\to \Tp(\Br(X_\kbar))$ is $\Tp(K_2)$. As above, we want to show that $\varprojlim_n K_{1,n}\to \Tp(K_2)$ is surjective. Combining the previous discussion and the fact that $\Br(k)$ is $p$-divisible by \cite[Thm. 1.3.7]{CS21}, we deduce that the two groups sit in the following diagram with exact rows
	
	\begin{equation*}\begin{tikzcd}
			0\arrow[r] &\Tp(\Br(k))\arrow[d,""]\arrow[r] &  \varprojlim_n K_{1,n}\arrow[d,""]\arrow[r] & H^1_{}(k,\Pic_{X/k}[p^n])\arrow[r]\arrow[d,""] & 0\\
			0\arrow[r]& \Tp(\Br(k)) \arrow[r]& \Tp(K_2)\arrow[r]& \Tp(H^1_{}(k,\Pic_{X/k}))\arrow[r]& 0.\
		\end{tikzcd}
	\end{equation*}

%The modules are free and of finite type by \cite[Chap. II, Prop. 5.9]{Ill79}.

For every $n>0$, the kernel of $H^1_{}(k,\Pic_{X/k}[p^n]) \to H^1_{}(k,\Pic_{X/k})[p^n]$ is $\Pic(X)/p^n$ and the groups $(\Pic(X)/p^n)_{n>0}$ form a Mittag--Leffler system. We deduce that the morphism $$\varprojlim_{n} H^1_{}(k,\Pic_{X/k}[p^n])\to \Tp(H^1_{}(k,\Pic_{X/k}))$$ is surjective. This implies that (\ref{2:e}) is exact, as we wanted.
\end{proof}



The proof of the following lemma was inspired by \cite[Thm. 2.1]{CS13}.

\begin{prop}\label{tricky:p}For an abelian variety $A$ over $k$, we have $H^2(A_\kbar,\Zp(1))^{k}[\tfrac{1}{p}]=(\varprojlim_{n}H^2(A_\kbar,\mupn)^{k})[\tfrac{1}{p}]$ and $\Tp(\Br(A_{\kbar}))^{k}[\tfrac{1}{p}]=\Vp(\Br(A_{\kbar})^{k}).$
\end{prop}

\begin{proof}
Thanks to \cite[Thm. 11]{Kat99}, there exists a proper smooth connected curve $C$ with a rational point and a morphism $C\to A$ such that $B:=\mathrm{Jac}(C)$ maps surjectively to $A$. By Poincaré's complete reducibility theorem, $B$ is isogenous to a product $A\times A'$ with $A'$ an abelian variety over $k$. Since $H^2(A_\kbar,\Zp(1))^{k}$ (resp. $\Br(A_{\kbar}))^k$) is a direct summand of $H^2(A_\kbar\times A'_\kbar,\Zp(1))^{k}$ (resp. $\Br(A_\kbar\times A'_{\kbar})^k$) and the property we want to prove is invariant by isogenies, it is then enough to prove the result for $B$. 
\spa

Let $K_{1,n}$ be the kernel of the morphism $H^2(B,\mupn)\to H^2(B_\kbar,\mupn)$. Since, by Lemma \ref{fppf-exte-scal:l}, the group $K_{1,n}$ is an extension of $H^1(k,\Pic_{B/k}[p^n])$ by $\Br(k)[p^n]$ and $\Pic_{C/k}[p^n]=\Pic_{B/k}[p^n]$, we deduce that $K_{1,n}=\ker(H^2(C,\mupn)\to H^2(C_\kbar,\mupn))$. By \cite[Rmq. 2.5.b]{Gro68}, the group $\Br(C_\kbar)$ vanishes, thus $H^2(C_{\kbar},\Zp(1))=\Zp$ and the morphism $H^2(C,\Zp(1))\to H^2(C_{\kbar},\Zp(1))$ is surjective because $C$ has a rational point. This implies that $R^1 \varprojlim_{n}K_{1,n}\to R^1 \varprojlim_{n}H^2(C,\mupn)$ is injective. Since $$R^1 \varprojlim_{n}K_{1,n}\to R^1 \varprojlim_{n}H^2(C,\mupn)$$ factors through $R^1 \varprojlim_{n}H^2(B,\mupn),$ we deduce that $R^1 \varprojlim_{n}K_{1,n}\to R^1 \varprojlim_{n}H^2(B,\mupn)$ is injective as well. Therefore, the morphism $H^2(B,\Zp(1))\to \varprojlim_{n}H^2(B_\kbar,\mupn)^{k}$ is surjective. Thanks to Proposition \ref{Kummer-k-inv:p}, this also implies that $\Tp(\Br(B_{\kbar}))^{k}\to\varprojlim_n(\Br(B_\kbar)[p^n])^k$ is surjective.

\spa

It remains to prove that $\mathrm{coker}(\ \varprojlim_n(\Br(B_\kbar)[p^n])^k\to \Tp(\Br(B_\kbar)^k)\ )$ has finite exponent. Combining Lemma \ref{fppf-exte-scal:l} and the fact that $\Br(C_\kbar)=0$, we deduce that there is an exact sequence $$0\to K_3'\to K_3\to \Br(B_\kbar)^k\to 0$$ with $K_3:=\ker(\Br(B)\to \Br(C))$ and $K_3':=\ker(H^1(k,\Pic_{B/k})\to H^1(k,\Pic_{C/k}))$. Since $\Pic^0_{C/k}=\Pic^0_{B/k}$, the group $K_3'$ is a subgroup of $H^1(k,\Pic_{B/k}/\Pic^0_{B/k})$, which is a finite group. In particular, $K_3'[p^\infty]$ is killed by $p^m$ for some $m>0$. Therefore, $\Br(B)[p^\infty]\to \Br(B_\kbar)^k[p^\infty]$ splits up to multiplication by $p^m$. We deduce that the image of $\varprojlim_n(\Br(B_\kbar)[p^n])^k\to \Tp(\Br(B_\kbar)^k)$ contains $p^m\Tp(\Br(B_\kbar)^k)$. This yields the desired result. 
\end{proof}


\section{Constructing a morphism}\label{cons-morp:s}
\label{cons-a-morp:s}
Let $A$ be an abelian variety over $k$. We want to complete the following natural square
	\begin{equation*}
	\begin{tikzcd}\Pic(A)^\wedge\arrow[r,"c_1"]\arrow[d] & H^2(A,\Zp(1))\arrow[d,dotted]\\
		\Hom(A,A^\vee)_{\Zp}\arrow[r] & \Hom_{p\textrm{-div}}(A[p^\infty],A^\vee[p^\infty]).
	\end{tikzcd}
\end{equation*}	

Over an algebraically closed field of characteristic $0$ such a commutative diagram is constructed in \cite[Lem. 2.6]{OSZ21}.

\subsection{}
	Consider the morphism $h_1:H^2(A,\mupn)\to H^2(A\times A,\mupn)$ which sends a class $\alpha$ to $m^*(\alpha)-\pi_1^*(\alpha)-\pi_2^*(\alpha)$, where $\pi_1$ and $\pi_2$ are the two projections of $A\times A$. This morphism has the property that for a line bundle $\calL$, its first Chern class $c_1(\calL)\in H^2(A,\mupn)$ is sent to $c_1(\Lambda(\calL))$, the first Chern class of the associated Mumford bundle $\Lambda(\calL)$. The Leray spectral sequence
\begin{equation}\label{lerayAxA:e}
	E^{i,j}_2=H^i(A,R^j\pi_{2*}\mupn)\Rightarrow H^{i+j}(A\times A,\mupn).
\end{equation}
induces a filtration $0\subseteq F^2 H^{2}(A\times A,\mupn) \subseteq F^1 H^{2}(A\times A,\mupn)\subseteq H^{2}(A\times A,\mupn) $.
\begin{lemm}\label{image:l}The image of $h_1$ lies in $F^1(H^{2}(A\times A,\mupn))$.
\end{lemm}
\begin{proof}
The spectral sequence (\ref{lerayAxA:e}) gives the exact sequence $$0\to F^1(H^{2}(A\times A,\mupn))\to H^{2}(A\times A,\mupn)\to E^{0,2}_\infty\to 0.$$ Therefore, it is enough to check that the composition $$H^2(A,\mupn)\xrightarrow{h_1} H^2(A\times A,\mupn)\to H^0(A,R^2\pi_{2*}\mupn)$$ is the $0$-morphism. By \cite{BO21} there exists a commutative algebraic group $G/k$ which represents $R^2q_{2*}\mupn$ on the big fppf site of $\Spec(k)$. This implies that $H^0(A,R^2\pi_{2*}\mupn)$ can be computed as $\Mor(A,G).$ The neutral component of $G$, denoted by $G^\circ$, is killed by the multiplication by $p^n$. Thus every homomorphism $A\to G^\circ$ is the $0$-morphism. We deduce that $\Mor(A,G)=\Mor(0_A,G)=H^0(k,R^2q_{*}\mupn)$. By \ref{limi:l}, the group $H^0(k,R^2q_{*}\mupn)$ is naturally a subgroup of $H^2(A_\kbar,\mupn)$ The induced morphism $$H^2(A\times A,\mupn)\to H^0(k,R^2q_{*}\mupn)\hookrightarrow H^2(A_\kbar,\mupn)$$ is induced by the pullback via $i_1:A=A\times 0_A \hookrightarrow A\times A$. By construction, we have that $i_1^*\circ h_1=i_1^*\circ m^*-i_1^*\circ \pi_1^*=0$. This concludes the proof.
\end{proof}

\subsection{}
Thanks to Lemma \ref{image:l}, we can define $$\bar{h}_1: H^2(A,\mupn)\to H^1(A,A^\vee[p^n])$$

as the composition of $h_1$ and the boundary morphism $F^1(H^2(A\times A,\mupn))\to H^1(A,R^1\pi_{2*}\mupn)=H^1(A,A^\vee[p^n])$.

\spa

We want to construct now a morphism $$h_2:H^1(A, A^\vee[p^n])\to\Hom(A[p^n],A^\vee[p^n]).$$ For an $A^\vee[p^n]$-torsor $P$ over $A$, let us associate a morphism $h_2(P): A[p^n]\to A^\vee[p^n]$. For every scheme $T/k$, a $T$-point of $A[p^n]$ corresponds by duality to a morphism $\tau: A_T^\vee[p^n]\to \mathbb{G}_{m,T}$. We define $h_2(P)(\tau)$ as $\tau_*(P)\in H^1(A_T,\mathbb{G}_{m,T})$. Finally, we define $$h: H^2(A,\mupn)\to \Hom(A[p^n],A^\vee[p^n])$$ as ${h_2}\circ\bar{h}_1$. We denote in the same way the induced morphism $H^2(A,\Zp(1))\to \Hom(A[p^\infty],A^\vee[p^\infty])$

\begin{prop}\label{comm-squa:p}
The square
		\begin{equation}
	\begin{tikzcd}
		\Pic(A)/p^n \arrow[d,""]\arrow[r,"c_1"] & H^2_{}(A,\mupn) \arrow[d,"h"]\\
		\Hom(A,A^\vee)/p^n \arrow[r]& \Hom(A[p^n],A^\vee[p^n]).\
	\end{tikzcd}
\end{equation}

is commutative.
\end{prop}
\begin{proof}
We have to show that for a line bundle $\calL\in \Pic(A)$, the homomorphism $$h(c_1(\calL))\in  \Hom(A[p^n],A^\vee[p^n])$$ is equal to $\varphi_{\calL}|_{A[p^n]}$. By construction, we have that $h_1(c_1(\calL))=c_1(\Lambda(\calL))$. Consider the Leray spectral sequence 
\begin{equation}\label{lerayAxAv:e}
E^{i,j}_2=H^i(A^\vee,R^j\pi_{2*}\mupn)\Rightarrow H^{i+j}(A\times A^\vee,\mupn).
\end{equation}


The morphism $A\times A\xrightarrow{\id_A\times \varphi_{\calL}} A\times A^\vee$ induces via pullback a morphism from (\ref{lerayAxAv:e}) to (\ref{lerayAxA:e}). If $\calP\in \Pic({A\times A^\vee})$ is the Poincaré bundle of $A$, we know that $\Lambda(\calL)=(\id_A\times \varphi_{\calL})^* \calP.$ Therefore, by functoriality, it is enough to show that the boundary morphism $F^1H^2(A\times A^\vee,\mupn)\to H^1(A^\vee,R^1\pi_{2*}\mupn)=H^1(A^\vee,A^\vee[p^n])$ sends $c_1(\calP)\in F^1H^2(A\times A^\vee,\mupn)$ to the $A^\vee[p^n]$-torsor $A^\vee\xrightarrow{\cdot p}A^\vee$. To prove this we look at the composition $H^1(A\times A^\vee,\Gm)\to H^0(A^\vee,R^1\pi_{2*}\Gm)\to H^1(A^\vee,R^1\pi_{2*}\mupn)$, where the first morphism is the natural one and the second one is the boundary morphism associated to the Kummer exact sequence for $\mupn$. The line bundle $\calP$ is first sent to the class associated to the tautological inclusion $A^\vee \hookrightarrow \Pic_{A/k}$, then it goes to the class of the tautological $A^\vee[p^n]$-torsor over $A^\vee$, as we wanted.
\end{proof}



\begin{prop}\label{imag:p}
The morphism $h:H^2(A_\kbar,\Zp(1))\to \Hom_{}(A_\kbar[p^\infty],A^\vee_\kbar[p^\infty])$ is an injective morphism with image $\Hom^\sym_{}(A_\kbar[p^\infty],A^\vee_\kbar[p^\infty])$, the group of homomorphisms which are fixed by the involution $\tau\mapsto \tau^\vee$.
\end{prop}

\begin{proof}	Write $W$ for the ring of Witt vectors of $\kbar$. By \cite[Rmq. II.5.10]{Ill79}, there is a canonical isomorphism $H^2(A_\kbar,\Zp(1))=H^2_{\crys}(A_\kbar/W)^{F=p}$. The morphism $H^2_{\crys}(A_\kbar/W)\to H^2_{\crys}(A_\kbar\times A_\kbar/W)$ which sends $\alpha$ to $m^*(\alpha)-\pi_1^*(\alpha)-\pi_2^*(\alpha)$ corresponds to the natural embedding $\bigwedge^2 H^1_{\crys}(A_\kbar/W)\hookrightarrow H^1_{\crys}(A_\kbar/W)^{\otimes 2}\subseteq H^2_{\crys}(A_\kbar\times A_\kbar/W)$, where the isomorphism $\bigwedge^2 H^1_{\crys}(A_\kbar/W)\simeq H^2_{\crys}(A_\kbar/W)$ is induced by the cup product and the inclusion $H^1_{\crys}(A_\kbar/W)^{\otimes 2}\subseteq H^2_{\crys}(A_\kbar\times A_\kbar/W)$ is the one produced by the K\"unneth formula. Since $H^1_{\crys}(A_\kbar/W)=H^1_{\crys}(A_\kbar^\vee/W)^\vee(-1)$, we have that $$(H^1_{\crys}(A_\kbar/W)^{\otimes 2})^{F=p}=\Hom_{\mathbf{F\textbf{-}Crys}}(H^1_{\crys}(A_\kbar^\vee/W),H^1_{\crys}(A_\kbar/W)),$$ where $\mathbf{F\textbf{-}Crys}$ is the category of $F$-crystals over $\kbar$. The $F$-crystals $H^1_{\crys}(A_\kbar/W)$ and $H^1_{\crys}(A_\kbar^\vee/W)$ are the contravariant crystalline Dieudonné modules of the $p$-divisible groups $A_\kbar[p^\infty]$ and $A^\vee_\kbar[p^\infty]$, thus we get $$(H^1_{\crys}(A_\kbar/W)^{\otimes 2})^{F=p}=\Hom(A_\kbar[p^\infty],A^\vee_\kbar[p^\infty]).$$
This yields the desired result.
\end{proof}
\section{Main results}\label{main-resu:s}

We are now ready to prove our main result, which is a flat version of the Tate conjecture for divisors of abelian varieties.

\begin{theo}\label{fppf-tate:t}
	Let $A$ be an abelian variety over a finitely generated field $k$ of characteristic $p$. The cycle class map $$\NS(A)_\Zp\xrightarrow{c_1} \varprojlim_n H^2(A_\kbar,\mupn)^k$$ is an isomorphism. In addition, $\Tp(\Br(A_\kbar)^k)=0$.
\end{theo}
\begin{proof} 
By Proposition \ref{Kummer-k-inv:p}, the cokernel of $\NS(A)_\Zp\xrightarrow{c_1} \varprojlim_n H^2(A_\kbar,\mupn)^k$ is torsion-free, thus it is harmless to invert $p$. The $\Zp$-module $\Hom(A[p^\infty],A^\vee[p^\infty])$ embeds into $\Hom(A_\kbar[p^\infty],A^\vee_\kbar[p^\infty])$, thus the morphism $$h:H^2(A,\Zp(1))\to \Hom(A[p^\infty],A^\vee[p^\infty])$$ induces a morphism $\tilde{h}:H^2(A_\kbar,\Z_p(1))^{k}\to \Hom_{p\textrm{-div}}(A[p^\infty],A^\vee[p^\infty])$. By Proposition \ref{imag:p}, we know that $\tilde{h}$ is injective and $\mathrm{im}(\tilde{h})$ is contained in $\Hom_{p\textrm{-div}}^\sym(A[p^\infty],A^\vee[p^\infty])$.
	In addition, by Proposition \ref{comm-squa:p}, we have the following commutative square	
	\begin{equation*}
		\begin{tikzcd}
			\NS(A)_{\Zp}\arrow[r,"c_1",hook]\arrow[d,hook] & H^2(A_\kbar,\Z_p(1))^{k}\arrow[d,"\tilde{h}",hook]\\
			\Hom(A,A^\vee)_{\Zp}\arrow[r] & \Hom_{p\textrm{-div}}(A[p^\infty],A^\vee[p^\infty]).
		\end{tikzcd}
	\end{equation*}
The lower arrow is an isomorphism by \cite[Thm. 2.6]{deJ98}. Since $\NS(A)_{\Q}=\Hom^\sym(A,A^\vee)_{\Q}$, we deduce that $\mathrm{im}(\tilde{h}\circ c_1)[\tfrac{1}{p}]=\Hom_{p\textrm{-div}}^\sym(A[p^\infty],A^\vee[p^\infty])[\tfrac{1}{p}]$. This implies that $c_1:\NS(A)_\Qp\to H^2(A_\kbar,\Z_p(1))^{k}[\tfrac{1}{p}]$ is surjective, thus by Proposition \ref{tricky:p}, we get that $c_1:\NS(A)_\Qp\to \break \varprojlim_n H^2(A_\kbar,\mupn)^k[\tfrac{1}{p}]$ is surjective, as we wanted. For the second part just note that by Proposition \ref{Kummer-k-inv:p}, the vector space $\Tp(\Br(A_\kbar))^k[\tfrac{1}{p}]$ vanishes and thanks to Proposition \ref{tricky:p} we deduce that $\Tp(\Br(A_\kbar)^k)=\Tp(\Br(A_\kbar)^k)[\tfrac{1}{p}]=0$.\end{proof}


%	\begin{equation*}
%		\begin{tikzcd}
%			\Pic(A)^\wedge\arrow[r]\arrow[d,hook] & \Hom^\sym_{p\textrm{-div}}(A[p^\infty],A^\vee[p^\infty])\arrow[d]\\
%			\Hom(A,A^\vee)_{\Zp}\arrow[r,"\sim"] & \Hom_{p\textrm{-div}}(A[p^\infty],A^\vee[p^\infty]).
%		\end{tikzcd}
%	\end{equation*}	







\begin{theo}\label{fini-Brau:t}
	Let $A$ be an abelian variety over a finitely generated field $k$ of characteristic $p>0$. The transcendental Brauer group $\Br(A_{\ks})^{k}$ has finite exponent. In particular, if the Witt vector cohomology group $H^2(A_{\kbar},W\calO_{A_\kbar})$ is a finite $W(\kbar)$-module, then $\Br(A_{\ks})^{k}$ is finite.
\end{theo}
\begin{proof}By \cite[Thm. 16.2.3]{CS21}, the group $\Br(A_{\ks})^{k}[\tfrac{1}{p}]$ is finite. Thanks to Corollary \ref{Br-inj:c}, the morphism $\Br(A_{\ks})\to \Br(A_{\kbar})$ is injective, which implies that $\Br(A_{\kbar})^{k}$ coincides with $\Br(A_{\ks})^{k}$. 
	Write $q: A_\kbar \to \Spec(\kbar)$ for the structural morphism. By \cite[Cor. 1.4]{BO21}, for every $n$ there exists an algebraic group $G_n$ representing $R^2q_{*}\mu_{p^n}$\footnote{One could use alternatively perfect groups rather than algebraic groups, as in \cite[Lem. 1.8]{Mil86}.}. Thanks to \cite[Prop. 3.1]{Mil86}, if $G_n^\circ$ is the neutral component of $G_n$, for $n$ big enough all the $G_n^\circ$ are isomorphic. This implies that there exists a finite exponent group $M$ and a natural number $a$, such that the group $\Br(A_\kbar)[p^\infty]$ is isomorphic to $(\Qp/\Zp)^{\oplus a}\oplus M$. If $H^2(A_{\kbar},W\calO_{A_\kbar})$ is a finite $W(\kbar)$-module, then $M$ is actually finite since each $G_n$ is finite by \cite[Cor. 12.5]{BO21}.  
	
	\spa
	
	 We are reduced to prove that $\Br(A_{\kbar})^{k}[p^\infty]$ has finite exponent. Suppose by contradiction that this is not the case. We deduce that $\Br(A_{\kbar})^{k}[p^\infty]\cap (\Qp/\Zp)^{\oplus a}$ has infinite exponent as well, which implies that $\Br(A_{\kbar})^{k}[p^\infty]$ contains a copy of $\Qp/\Zp$. On the other hand, by Theorem \ref{fppf-tate:t}, the group $\Tp(\Br(A_{\kbar})^{k})$ vanishes, which leads to a contradiction.
\end{proof}

\begin{coro}\label{galo-fixe:c}
The group $\Br(A_{\ks})^{\Gamma_k}$ has finite exponent.
\end{coro}
\begin{proof}
This follows from Theorem \ref{fini-Brau:t} by arguing as in the proof of \cite[Thm. 2.1]{CS13}.
\end{proof}







\section{Specialisation of the Néron--Severi group}
\label{spec:s}

\subsection{}
We want to start this section with an explicative example. Let $\calE\to X$ be a non-isotrivial family of ordinary elliptic curves, where $X$ is a connected normal scheme of finite type over $\Fp$. Let $\calA$ be the product $\calE\times \calE$. We denote by $E$ and $A$ the generic fibres over the generic point $\Spec(k)\hookrightarrow X$. The Kummer exact sequence induces the exact sequence 
	$$0\to\mathrm{NS}(A_{\bar{k}})_{\Z_p} \to H^2_{}(A_{\bar{k}},\Z_p(1))\to \Tp(\mathrm{Br}(A_{\bar{k}}))\to 0.$$

The group $\mathrm{NS}(A_{\bar{k}})_{\Z_p}$ is of rank $2+\rk_\Z(\End_k(E))=3$, while the rank of $H^2_{}(A_{\bar{k}},\Z_p(1))$ is $2+\rk_\Zp(\End(E_\kbar[p^\infty]))=4$. This implies that $\Tp(\mathrm{Br}(A_{\bar{k}}))$ is of rank $1$. The endomorphisms of $E_{\kbar}[p^\infty]$ are all defined over $\ki$, which implies that the action of $\Gamma_k$ on $H^2_{}(A_{\bar{k}},\Z_p(1))$ is trivial. In particular, the morphism $\mathrm{NS}(A)_{\Z_p} \to H^2_{}(A_{\bar{k}},\Z_p(1))^{\Gamma_k}$ is not surjective and the cokernel $\Tp(\mathrm{Br}(A_{\bar{k}}))^{\Gamma_k}$ is isomorphic to $\Zp$. This shows that the Galois action on $H^2_{}(A_{\bar{k}},\Z_p(1))$ is not enough to detect what classes are $\Zp$-linear combinations of algebraic cycles. There is an additional obstruction to descend cohomology classes through the purely inseparable extension $\ki/k$.

\spa

This extra purely inseparable obstruction gives an explanation of the failure of surjectivity of the specialisation morphisms for the Néron--Severi group. In the example, if $\Spec(\kappa)\hookrightarrow X$ is a closed point, we have that $\NS(\calA_{{\kappa}})$ is of rank $4$ because $\End_\kappa(\calE_\kappa)$ is of rank $2$ (there is an extra Frobenius endomorphism). Thus, the specialisation map $\NS(A)\hookrightarrow \NS(\calA_\kappa)$ is never surjective even if the rank of $H^2_{}(A_{\bar{\kappa}},\Z_p(1))$ is $4$ as the generic geometric fibre. Since $\kappa$ is perfect, the extra obstruction coming from the purely inseparable extension $\ki/k$ does not affect $H^2_{}(A_{\bar{\kappa}},\Z_p(1))$. In general, we prove the following theorem.

\begin{theo}\label{NS:t}
	Let $X$ be a connected normal scheme of finite type over $\Fp$ with generic point $\eta=\Spec(k)$ and let $f:\calA\to X$ be an abelian scheme over $X$ with constant slopes. For every closed point $x=\Spec(\kappa)$ of $X$ we have $$\rk_\Z(\NS(\calA_{\bar{x}})^{\Gamma_\kappa})-\rk_\Z(\NS(\calA_{\bar{\eta}})^{\Gamma_k})\geq \rk_\Zp ( \Tp(\Br(\calA_{\bar{\eta}}))^{\Gamma_k}).$$
\end{theo}

For this purpose we first need the following result.

\begin{prop}\label{lisse-sheaf:p}
	The functor $\calF$ which sends $T\in (X^{\mathrm{perf}})_{\mathrm{pro}\et}$\footnote{We denote by $(-)_{\mathrm{pro}\et}$ the pro-étale site of a scheme, as defined in \cite{BS15}.} to $\Hom^{\sym}(\calA_T[p^\infty],\calA^\vee_T[p^\infty])[\tfrac{1}{p}]$ is a semi-simple $\Qp$-local system.
\end{prop}
\begin{proof} To prove that $\calF$ is a $\Qp$-local system one can imitate the proof of \cite[Prop. 4.3.13]{CS17}. For the semi-simplicity, we can shrink $X$ and assume it smooth. Write $A$ for the generic fibre of $\calA$. As we have seen in Proposition \ref{imag:p}, there are canonical isomorphisms $$H^2(A_{\kbar},\Zp(1))=\Hom^\sym(A_{\kbar}[p^\infty],A^\vee_{\kbar}[p^\infty])	=\Hom^\sym_{\mathbf{F\textbf{-}Crys}}(H^1_{\crys}(A_{\kbar}^\vee/W(\kbar)),H^1_{\crys}(A_{\kbar}/W(\kbar))),$$
	
	where we denote again with the superscript $^\sym$ the group of homeomorphisms which are fixed by the involution $\tau\mapsto \tau^\vee$. The second identification descends to $\ki$, since by Dieudonné's theory we also have $$\Hom^\sym(A_{\ki}[p^\infty],A^\vee_{\ki}[p^\infty])=\Hom^\sym_{\mathbf{F\textbf{-}Crys}}(H^1_{\crys}(A_{\ki}^\vee/W(\ki)),H^1_{\crys}(A_{\ki}/W(\ki))).$$
	
 At the same time, thanks to the crystalline K\"unneth formula, we have $$\Hom^\sym_{\mathbf{F\textbf{-}Crys}}(H^1_{\crys}(A_{\ki}^\vee/W(\ki)),H^1_{\crys}(A_{\ki}/W(\ki)))=H^2_{\crys}(A_\ki/W(\ki))^{F=p}.$$We write $N$ for the unit-root part of the $F$-crystal $H^2_{\crys}(A_\ki/W(\ki))(1)$ over $\Spec(\ki)$.
	
	\spa
	
	Thanks to \cite[Prop. 4.1.1]{Kat73}, there exists a natural equivalence between unit-root $F$-crystals over $\Spec(\ki)$ and finite free $\Zp$-modules with a continuous representation of $\Gamma_{\ki}$. The equivalence sends an $F$-crystal $M$ to the pro-étale torsor $T/\ki$ such  that $$T(k'):=\{\textrm{Frobenius-fixed bases of $M\otimes_{W(\ki)} W(k')$} \}$$ for every algebraic extension $k'/\ki$. Thanks to the previous discussion and the functoriality of the Dieudonné modules, we deduce that the Galois representation of $\Gamma_{\ki}$ on $H^2(A_{\kbar},\Zp(1))$ corresponds to Katz's pro-étale torsor associated to $N$.
	
	\spa
	
By \cite[Theorem 2.1]{Cre87}, there is a similar equivalence between unit-root $F$-isocrystals over $X$ and finite dimensional $\Qp$-local systems. By construction, Crew and Katz's correspondences are compatible, in the sense that they agree after pulling back the objects via $\eta^\mathrm{perf}\to X$. Write $\calM$ for the $F$-isocrystal $(R^2f_{\crys*}\calO_{\calA,\crys})(1)$ and $\calN$ for the quotient $\calM^{\geq 0}/\calM^{<0}$, where $\calM^{\geq 0}$ (resp. $\calM^{<0}$) is the subobject of $\calM$ of slopes $\geq 0$ (resp. $<0$). The $F$-crystal $N$, when we invert $p$, is the same as the fibre of $\calN$ at $\eta^{\mathrm{perf}}$. By \cite[Thm. 1.1.2]{D'Ad20}, the $F$-isocrystal $\calN$ is semi-simple. This yields the desired result.\end{proof}




\subsection{}\textit{Proof of Theorem \ref{NS:t}. }We look at the exact sequence
$$0\to\NS(A_{\kbar})_{\Qp}\to H^2(A_{\kbar},\Qp(1))\to \Tp(\Br(A_{\kbar}))[\tfrac{1}{p}]\to 0.$$

Thanks to Proposition \ref{lisse-sheaf:p}, the operation of taking Galois fixed points is exact. We get the exact sequence
$$0\to\NS(A_{\kbar})_{\Qp}^{\Gamma_k}\to H^2(A_{\kbar},\Qp(1))^{\Gamma_k}\to \Tp(\Br(A_{\kbar}))[\tfrac{1}{p}]^{\Gamma_k}\to 0.$$
Looking at the ranks we deduce the following equality
\begin{equation}\label{eq:e}
\rk_\Zp(H^2(A_{\kbar},\Zp(1))^{\Gamma_k})=\rk_\Z(\NS(A_{\kbar})^{\Gamma_k})+\rk_\Zp(\Tp(\Br(A_{\kbar}))^{\Gamma_k}).
\end{equation}
By Proposition \ref{lisse-sheaf:p}, the action of $\Gamma_k$ on $H^2(A_{\kbar},\Qp(1))=\Hom^\sym(A_\kbar[p^\infty],A^\vee_\kbar[p^\infty])[\tfrac{1}{p}]$ factors through the étale fundamental group of $X$ associated to $\bar{\eta}$, denoted by $\pi_1^{\et}(X,\bar{\eta})$. If $\kappa$ is the residue field of $x$, the inclusion $x\hookrightarrow X$ induces then an action of $\Gamma_\kappa$ on $H^2(A_{\kbar},\Qp(1))$ which corresponds, up to conjugation, to the action of $\Gamma_\kappa$ on $H^2(A_{\bar{\kappa}},\Qp(1))$. Therefore, by Corollary \ref{galo-fixe:c}, we get $\NS(\calA_{\bar{x}})_\Qp=H^2(A_{\kbar},\Qp(1))^{\Gamma_\kappa}$. Since $H^2(A_{\kbar},\Qp(1))^{\Gamma_k}$ is a subspace of $H^2(A_{\kbar},\Qp(1))^{\Gamma_\kappa}$ we show that
\begin{equation}\label{ineq:e}
	\rk_\Z(\NS(\calA_{\bar{x}})^{\Gamma_\kappa})=\rk_\Zp(H^2(A_{\kbar},\Zp(1))^{\Gamma_\kappa})\geq \rk_\Zp(H^2(A_{\kbar},\Zp(1))^{\Gamma_k}).
\end{equation}

Combining (\ref{eq:e}) and (\ref{ineq:e}) we deduce the desired result.\qed

\spa

We want to conclude this section with some more general examples of abelian varieties such that $ \Tp(\Br(A_{\kbar}))^{\Gamma_k}\neq 0$. These abelian varieties are all counterexamples to the conjecture in \cite[§7.3.1]{Ulm14} when $\ell=p$. Let $A$ be an abelian variety which splits as a product $B\times B$ with $B$ an abelian variety over $k$.
\begin{prop}\label{counter:p}
		Let $B$ be an abelian variety over $k$. There is a natural exact sequence $$0\to \Hom(B,B^\vee)_\Zp\to \Hom(B_\kbar[p^\infty],B_\kbar^\vee[p^\infty])^{\Gamma_k}\to \Tp(\Br(A_{\kbar}))^{\Gamma_k},$$ where $A:=B\times B$.
			
	
\end{prop}


\begin{proof}
	We consider the exact sequence $$0\to \NS(A_{\kbar})^{\Gamma_k}_\Zp\to H^2(A_{\kbar}, \Zp(1))^{\Gamma_k}\to \Tp(\Br(A_{\kbar}))^{\Gamma_k}.$$ Arguing as in the proof of Proposition \ref{imag:p}, the $\Zp$-module $\Hom(B_\kbar[p^\infty],B_\kbar^\vee[p^\infty])^{\Gamma_k}$ is naturally a direct summand of $H^2(A_{\kbar}, \Zp(1))^{\Gamma_k}$. Its preimage in $\NS(A_{\kbar})_\Zp^{\Gamma_k}$ corresponds to the $\Zp$-module $\Hom(B_\kbar,B^\vee_\kbar)_\Zp^{\Gamma_k}=\Hom(B,B^\vee)_\Zp.$ This concludes the proof.
\end{proof}



\begin{coro}\label{count:c}
	If $\End(B)=\Z$, then $\Tp(\Br(A_{\kbar}))^{\Gamma_k}\neq 0$.
\end{coro}

\begin{proof}
By the assumption, $\Hom(B,B^\vee)_\Zp$ is a $\Zp$-module of rank $1$. Therefore, by Proposition \ref{counter:p}, it is enough to prove that the rank of $\Hom(B_\kbar[p^\infty],B_\kbar^\vee[p^\infty])^{\Gamma_k}$ is greater than $1$. Since $\End(B)=\Z$, the abelian variety $B$ is not supersingular, so that the $p$-divisible group $B[p^\infty]$ admits at least two slopes. By the Dieudonné--Manin classification, this implies that $B_{\ki}[p^\infty]$ is isogenous to a direct sum $\calG_1\oplus \calG_2$ of non-zero $p$-divisible groups over $\ki$. Since $\End(\calG_1)[\tfrac{1}{p}]\oplus \End(\calG_2)[\tfrac{1}{p}]$ embeds into $\End(B_{\ki}[p^\infty])[\tfrac{1}{p}]\simeq \Hom(B_\kbar[p^\infty],B_\kbar^\vee[p^\infty])[\tfrac{1}{p}]^{\Gamma_k}$, we deduce that $\rk_\Zp(\Hom(B_\kbar[p^\infty],B_\kbar^\vee[p^\infty])^{\Gamma_k})>1$, as we wanted.
\end{proof}
%\begin{rema}
%Looking at monodromy groups, using the parabolicity conjecture the inequality is
%$$V^{F_x}-V^G\geq V^{L}-V^G$$

%where $G$ is the overconvergent monodromy, $L$ is the commutator of the slope cocharacter and $F_x\in G$ is the Frobenius at $x$.
%\end{rema}


\bibliographystyle{ams-alpha}
\begin{thebibliography}{ABCD1}
\bibitem[Amb18]{Amb18} A.~Ambrosi, \textit{Specialization of Néron-Severi groups in positive characteristic}, \href{https://arxiv.org/abs/1810.06481
}{\url{arXiv:1810.06481}} (2018), to appear in Ann. Sci. École Norm. Sup.


	
	
	\bibitem[AD18]{AD18} A.~Ambrosi and M.~D'Addezio, \emph{Maximal tori of monodromy groups of $F$-isocrystals and an application to abelian varieties}, \href{https://arxiv.org/abs/1811.08423}{\url{arXiv:1811.08423}} (2018).
	
	\bibitem[And96]{And96}Y. André, Pour une théorie inconditionnelle des motifs, \textit{Inst. Hautes Etudes Sci. Publ. Math.} 
	\textbf{83} (1996), 5--49.
	\bibitem[AM77]{AM77} M. Artin and B. Mazur, Formal groups arising from algebraic varieties,
	\textit{Ann. Sci. École Norm. Sup.} \textbf{10} (1977), 87--132.
		\bibitem[BO21]{BO21} D.~Bragg and M. Olsson, \emph{Representability of cohomology of finite flat abelian group schemes}, \href{https://arxiv.org/abs/2107.11492}{\url{arXiv:2107.11492}} (2021).
		
		\bibitem[BS15]{BS15}B.~Bhatt and P. Scholze, The pro-\'etale topology for schemes, \textit{Ast\'erisque} \textbf{369} (2015), 99--201.

%	\bibitem[CHT21]{CHT21}A. Cadoret, C. Y. Hui, and A. Tamagawa, \textit{$\Ql$- versus $\F_\ell$-coefficients in the Grothendieck-Serre-Tate conjectures}, preprint (2021).

\bibitem[CS17]{CS17}A. Caraiani and P. Scholze, On the generic part of the cohomology of compact unitary Shimura varieties, \textit{Ann. of Math.} \textbf{186} (2017), 649--766.
	
\bibitem[Chr18]{Chr18}	A. Christensen, \textit{Specialization of Néron-Severi groups in characteristic $p$}, \href{https://arxiv.org/abs/1810.06550}{\url{arXiv:1810.06550}} (2018).
	
	
	\bibitem[CS13]{CS13} J.--L. Colliot-Thélène and A. N. Skorobogatov, Descente galoisienne sur le groupe de Brauer, \emph{J. reine angew. Math.} \textbf{682} (2013), 141--165.
	
	\bibitem[CS21]{CS21} J.--L. Colliot-Thélène and A. N. Skorobogatov, \textit{The Brauer--Grothendieck group}, Springer, 2021.
	
	\bibitem[Cre87]{Cre87}R. Crew, $F$-isocrystals and $p$-adic representations, \textit{Proc. Symp. Pure Math.} \textbf{46} (1987), 111--138.
	
	\bibitem[D'Ad20]{D'Ad20}M. D'Addezio, \textit{Parabolicity conjecture of $F$-isocrystals}, \href{https://arxiv.org/abs/2012.12879}{\url{arXiv:2012.12879}} (2020).
	\bibitem[deJ98]{deJ98}A.~J.~de Jong, Homomorphisms of Barsotti--Tate groups and crystals in positive characteristic, \emph{Invent. Math.} \textbf{134} (1998), 301--333.
	%\bibitem[Gab01]{Gab01}O.~Gabber, On space filling curves and Albanese varieties, \textit{Geometric \& Functional Analysis GAFA} \textbf{11} (2001), 1192--1200.
	\bibitem[Gro68]{Gro68} A.~Grothendieck, Le groupe de Brauer III: Exemples et compléments, in \textit{Dix Exposés sur la Cohomologie des Schémas}, North-Holland, 1968, 88--188.
	

	\bibitem[Ill79]{Ill79}L.~Illusie, Complexe de de Rham-Witt et cohomologie cristalline, \textit{Ann. Sci. École Norm. Sup.} \textbf{12} (1979), 501--661.

    \bibitem[Ill83]{Ill83}
	L. Illusie, Finiteness, duality and K\"unneth theorems in the cohomology of the de Rham-Witt complex, in \textit{Algebraic Geometry, Proc. conf. Tokyo/Kyoto}, Springer, 1983, 20--72.
	
		\bibitem[Kat73]{Kat73}N. M. Katz, P-adic properties of modular schemes and modular forms, in \textit{Modular functions of one variable III}, Springer, 1973, 69--190.
	
	\bibitem[Kat99]{Kat99}
	N.~M.~Katz, Space filling curves over finite fields, \emph{Math. Res. Lett.} \textbf{6} (1999), 613--624.
	
    \bibitem[MP12]{MP12}D. Maulik and B. Poonen, Néron--Severi groups under specialization, \textit{Duke Math. J.} \textbf{161} (2012), 2167--2206.
    
    
	
	\bibitem[Mil86]{Mil86} J. Milne, Values of zeta functions of varieties over finite fields, \textit{Amer.
	J. Math.} \textbf{108} (1986), 297--360.
\bibitem[OSZ21]{OSZ21}M. Orr, A. N.~Skorobogatov, and Y. Zarhin, On uniformity conjectures for abelian varieties and K3 surfaces,\textit{ Amer. J. Math.} \textbf{143} (2021), 1665--1702. 
	
	\bibitem[SZ08]{SZ08} A. N.~Skorobogatov and Y. G. Zarhin, A finiteness theorem for the Brauer group of abelian varieties and K3 surfaces, \emph{J. Algebraic Geom.} \textbf{17} (2008), 481--502.
	
 \bibitem[SP22]{SP22}The Stacks Project Authors, \textit{Stacks Project}, \url{http://stacks.math.columbia.edu} (2022).
	
	\bibitem[Tat66]{Tat66} J. Tate, Endomorphisms of abelian varieties over finite fields, \textit{Invent. Math}. \textbf{2} (1966), 134--144.
	\bibitem[Tat94]{Tat94} J. Tate, Conjectures on algebraic cycles in $\ell$-adic cohomology, in \textit{Motives (Seattle, WA, 1991), Proc. Sympos. Pure
	Math.} \textbf{55}, Amer. Math. Soc., 1994, 71--83.
	\bibitem[Ulm14]{Ulm14}D. Ulmer, Curves and Jacobians over function fields, in \textit{Arithmetic geometry over global function fields}, Adv. Courses Math. CRM, Birkh\"auser-Springer, 2014, 283--337.
\end{thebibliography}

\end{document}
